\chapter{Desenvolvimento}
\label{cap:desenvolvimento}
Introdução do capitulo pendente 


\section{RISC-V}
\label{sec:risc-v}

As intrunçoes do conjunto RV32I são organizado dentro dos quattros formatos basicos exibidos na Figura \ref{fig:formatoInstrucaoRV32I_2}, porem para alguns cassos esses formatos puros não são suficiente para conteplar todas posibilidades, ou não são o mais eficiente, por causa disso exite algumas variçoes,  

\begin{figure}[!htb]%% Ambiente figure
     \centering % Centraliza a imagem no documento (recomendado)
     \caption{Formato das instruções da RV32I}% Legenda da sua figura
     \label{fig:formatoInstrucaoRV32I_2}%rotulo
     \includegraphics[scale=0.54]{figuras/instrucaoRISC-V.pdf}
     \fonte{Adaptado de \cite{riscv_manual_2026}}%% Fonte adaptada
\end{figure}

Dentro do conjunto RV32I as instruções são dividas em cinco categoraias: Inteiros, Transferência de controle, Carga e Descarga da memoria, Ordenação da memoria e Chamada de ambiente \cite{riscv_manual_2026}.

O conjunto dos inteiros é composto por 21 intruçoes e é o resposalvel por realizar as operaçoes matemticas soma e subtração, e operaçoes logicas de comapração, deslocamento, AND, OR ou XOR. A operação pode ser realizada tanto entre dois numeros vindo de registradores, quanto entre um registrador e um valor constante vindo diretemente da intrução que é nomeado de imediato.

O conjunto RV32I possui originalmente 40 instruções em sua especificação base; contudo, implementações simples como a proposta neste trabalho podem reduzir para 38 intruçoes. A primeira redução ocorre nas duas únicas instruções do tipo SYSTEM: a instrução \texttt{ECALL}, responsável por solicitar chamadas de sistema (\textit{system calls}) ao sistema operacional, e a instrução \texttt{EBREAK}, utilizada por depuradores para pausar a execução do programa e devolver o controle ao ambiente de desenvolvimento. Em um núcleo de ciclo único sem suporte a sistema operacional ou depuração em tempo real, ambas as instruções podem ser unificadas em uma rotina simplificada de exceção (\textit{trap}), disparando um sinal de interrupção ou paralisando a execução do circuito de forma controlada \cite{riscv_manual_2026}. 


Outra instrução que pode ser simplificada em implementações de escopo reduzido é a \texttt{FENCE}. Essa instrução é responsável por garantir o ordenamento na execução de memória em sistemas multicore ou em microarquiteturas avançadas com execução fora de ordem (\textit{out-of-order}). Como o processador projetado neste trabalho é de núcleo único e possui execução estritamente sequencial em ciclo único, a ordenação das operações de memória é garantida nativamente pela própria estrutura física. Dessa forma, a instrução \texttt{FENCE} torna-se redundante, podendo ser tratada funcionalmente como uma instrução NOP (do inglês \textit{No Operation}), isto é, uma operação nula que não altera o estado dos registradores nem da memória \cite{riscv_manual_2026}.


\section{Caminho de dados}
\label{sec:desenvolviemnto}

O desenvolvimento do hardware será estruturado de forma estritamente modular, dividindo a arquitetura do processador em blocos funcionais discretos e independentes, o que viabiliza a reutilização desses módulos em projetos futuros. Adota-se o modelo clássico de caminho de dados (\textit{datapath}) em ciclo único, no qual o fluxo de processamento de uma instrução ocorre integralmente dentro de um único período de \textit{clock}. A integração dessas unidades será realizada de forma hierárquica por meio de um módulo de topo (\textit{top-level}), o qual gerencia as interconexões de sinais internos, facilitando tanto as simulações comportamentais quanto a síntese física na plataforma Nexys A7.


A estrutura interna do processador será segmentada em cinco blocos principais: Unidade de Busca, composta pelo Contador de Programa (PC, do inglês \textit{Program Counter}) e por um somador dedicado; Unidade de Decodificação, constituída pelo circuito decodificador da instrução e pelo Banco de Registradores; Unidade de Execução (ALU); Unidade de Memória de dados; e Unidade de Escrita (\textit{Write-Back}). A interação simplificada entre esses blocos pode ser visualizada na Figura \ref{fig:comunicacaoBloco}. Essa partição do sistema foi projetada para garantir o isolamento de falhas durante os testes individuais e, estrategicamente, para facilitar a conversão futura da microarquitetura para o modelo segmentado em \textit{pipeline}, tomando como referência o caminho de dados estabelecido por \citeonline{hennessy2017organizacao} para um núcleo com suporte a sete instruções, \textit{ load doubleword}, \textit{store doubleword }, \textit{add}, \textit{sub}, \textit{and}, \textit{or} e \textit{beq} é ilustrado na Figura \ref{fig:caminhoDadosRiscV}.


\begin{figure}[!htb]
     \centering 
     \caption{Comunicação entre os blocos do processador}
     \label{fig:comunicacaoBloco}
     \includegraphics[scale=0.5]{figuras/FluxoProcessador.pdf}
     \fonte{Autoria própria (2026)}
     
\end{figure}

\begin{figure}[!htb]
     \centering 
     \caption{Caminho de dados de uma implementação RISC-V de ciclo único}
     \label{fig:caminhoDadosRiscV}
     \includegraphics[scale=1]{figuras/caminhoDadosRiscV.pdf}
     \fonte{Adaptado de \citeonline{hennessy2017organizacao}}
    
\end{figure}



Para a validação de cada bloco lógico, será desenvolvido um ambiente de teste isolado (\textit{testbench}), permitindo a simulação comportamental completa do circuito e a verificação de suas funcionalidades dentro da própria suíte AMD Vivado. Cada módulo será testado individualmente e, somente após a constatação de sua estabilidade operacional, proceder-se-á com a integração hierárquica, aplicando novos estímulos de simulação a cada etapa de junção do hardware.

\subsection{ALU}


\subsubsection{Somador}

Um dos metodos mais visto para somadores na literatura é utilizando um circuito chamado de somador completo, que consite em uma pequeno circuito resposavel por realizar a soma entre dois bits de um numero, assim como realizamos a soma normalmete esse circuito soma cada digito do numero separado e se "sobrar" passa adiante essa valor para somar no proximo digito, por exemplo para realizar a soma: 

\begin{equation}
    \opadd[carrystack=true]{191}{874}
    \label{eq:soma_colunas}
\end{equation}


Em binario a podemso fazer o mesmo, e essa é a logica do soamdor compelto, soamr dois bits, e o "vai um" se existir, assim o circuito fica com 3 entradas, os dois bits a ser soamdo, o bit de vem um, conhecido com o \textit{carry-in}, e 2 saidas, o resultado da som e o bit de vai um conhecido como \textit{carry-out}, esse circuito pode ser facilemte contruido quando anlisado a tabela verdade das daus saidas

% Tabela para a Saída do Bit de Soma (S)
\begin{table}[htb]
\centering
\caption{Tabela verdade do somador completo para a saída do bit de soma ($S$).}
\label{tab:fa_soma}
\begin{tabular}{cccc}
\toprule
\textbf{A} & \textbf{B} & \textbf{$C_{in}$} & \textbf{Soma ($S$)} \\
\midrule
0 & 0 & 0 & 0 \\
0 & 0 & 1 & 1 \\
0 & 1 & 0 & 1 \\
0 & 1 & 1 & 0 \\
1 & 0 & 0 & 1 \\
1 & 0 & 1 & 0 \\
1 & 1 & 0 & 0 \\
1 & 1 & 1 & 1 \\
\bottomrule
\end{tabular}
\fonte{Elaborado pelo autor}
\end{table}

% Tabela para a Saída do Carry-Out (Cout)
\begin{table}[htb]
\centering
\caption{Tabela verdade do somador completo para o transporte de saída ($C_{out}$).}
\label{tab:fa_cout}
\begin{tabular}{cccc}
\toprule
\textbf{A} & \textbf{B} & \textbf{$C_{in}$} & \textbf{$C_{out}$} \\
\midrule
0 & 0 & 0 & 0 \\
0 & 0 & 1 & 0 \\
0 & 1 & 0 & 0 \\
0 & 1 & 1 & 1 \\
1 & 0 & 0 & 0 \\
1 & 0 & 1 & 1 \\
1 & 1 & 0 & 1 \\
1 & 1 & 1 & 1 \\
\bottomrule
\end{tabular}
\fonte{Elaborado pelo autor}
\end{table}

resultando nas equaçoes booleanas:

$$S = A \oplus B \oplus C_{in}$$
e
$$C_{out} = (A \cdot B) + (C_{in} \cdot (A \oplus B))$$

Uma subtração nada mais é que a soma entre um numero positivo e um nuemro negativo, portanto para subtarir dois nuemos basta que trasnfomemos esse nuemro em negativo utilizando a o metodo para complemteto de dois, portanto podemos utilizar uma inversora e enviar um sinal de carry-in 1

tendo as equaçoes um método de criar um soamdor seria apenas criar um somador dentro do verilog utilizando a parte logica da FPGA, porem se feito assim o caminho critico aba ficando muito grande devido a progração dos bits, para resolver isso e melhorar esse efeito de propragação de bits as FPGA possui unidade propria de alta velocidade feita para acelerar o transporte de bit chamada carry4

explciar sobre o carry4

Para utilizar os blocos carry pode ser feito utilizando operadores "+" ou "-" que ja intanciam e adicionan toda a logica nescessaria, ou intanciando eles manualmete junto da logica nescessaria. Foi realizado um teste com um progrma simples que realiza a leitura dos 16 shitchs da placa, divide ele em dois nuemos de 8 bits, send A os oito primerios bits, e B os bits mais significativo, e o resultado seria mostrado nos 16 leds, novamenet dividos em dois pares de 8 bits conforme anteriormente, para conseguir medir o tempo dos dois circuitos foi nescessario colocar um banco de registradores na entrda dos shitchs e um banco de registradores n saida dos leds.

Foi desenvolvido dois modulos, o primeiro contendo um soamdor com a feramente automatica da Listing \ref{lst:sub_auto} e o segundo da Lissting \ref{lst:sub_manula} onde alem de estanciar os dois blocos carry4 nescessario foi precisso relizar tanto a operação  de inverter o op\_b e realizar a XOR com op\_a. 


% O trecho de código mencionado no texto não foi incluído.

Obserevando os diagrmas gerado pela feramente de schematci gerado pelo virado pós-implementação é posivel observar algumas pequenas defirenças, na Figura \ref{fig:sub_auto} do circuito gerado automatico é possível observar que foi criado oito LUT e tres carry4, ja na Figura \ref{fig:sub_manual} um bloco de carry4 foi substituido por uma LUT de 1 bit.
\begin{figure}[!htb]
     \centering 
     \caption{Diagrama do circuito resultante do circuito de subtração automático}
     \label{fig:sub_auto}
     \includegraphics[scale=1]{figuras/sub_auto.pdf}
     \fonte{Autoria própria (2026)}
    
\end{figure}

\begin{figure}[!htb]
     \centering 
     \caption{Diagrama do circuito resultante do circuito de subtração manual}
     \label{fig:sub_manual}
     \includegraphics[scale=1]{figuras/sub_manual.pdf}
     \fonte{Autoria própria (2026)}
    
\end{figure}

Somente observar o diagrama diz nada sobre efeiciencia, para realmete conseguir definir qual o melhor é preciso fazer uma comparação entre os tempo dos caminhos criticos, para conseguir determinar a distancia entre os caminhos criticos é preciso colocar um registrador em cada entrada e saida, parra assim a feramenta de analise de timing do proprio software vivado conseguir calcular o caminho critico. Portanto foi feito o código X que intancia os registradores  e os dois modulos.


\begin{lstlisting}[language=Verilog, caption={Modulo em SystemVerilog do Subtração para comaprativo de desempenho manula}, label={lst:sub_manula}]
module top_benchmark (
    input  logic        clk,
    input  logic [15:0] switch,
    output logic [8:0]  led_auto,
    output logic [8:0]  led_manual
);
    // Registradores de Entrada
    logic [15:0] switch_reg;
    always_ff @(posedge clk) switch_reg <= switch;

    // Sinais dos submódulos
    logic [8:0] res_auto, res_manual;

    // Instanciação
    sub_auto   u_auto   (.switch(switch_reg), .clk(clk), .led(res_auto));
    sub_manual u_manual (.switch(switch_reg), .clk(clk), .led(res_manual));

    // Registradores de Saída
    always_ff @(posedge clk) begin
        led_auto   <= res_auto;
        led_manual <= res_manual;
    end
endmodule
\end{lstlisting}

Com o circuito montado é executado na ferameta de compilaççã odo código até a etapada de implemetmnação, assim permitidno utilizar a analise de design de implemtnação, dentro desse relatório o inteeresante é analise do \textit{Setup timing} dos caminhos de circuito interno, o valor desses realtorio é o tempo nescessario para que uma informação saia do registradores internos, seja processado e chege no registrador de saida, com estabilidade suficiente para garantir que o valor que será salvo é realmet o valor coreto e não algum erro causando por transição de sinal.

A Tabela \ref{tab:ComTming} contem o caminho criticos de cada circuito, é posivel tirar uma conclução finalemte, a coluna level indica a quantidade de blocos logicos que existee no caminho critico, quantos mais blocos o caminho tende a ser mais lento, e o total delay é o tempo total relamte do coircuito em nano segundos. É evidente que o caminho automatico obteve u mdesempenho  muito superior ao cominho manual, gastando aproxiamndamente 1ns a menos.

Como ficou evidente o processo onde  aferamte monta o circuito de soma automaticamente se parece ser  o mais eficiente pelo menos comaprando com uma simples 

\begin{table}[!htb]
\caption{Comparativo de timign entre o metodo de soma\label{tab:ComTming}}
    \centering
    \begin{tabular}{ccc}
         & Level & Total Delay (ns) \\
        Manual & 4 & 3,221 \\
        Atuo & 3 & 2,649 \\
    \end{tabular}
    \caption{Caption}
    \label{tab:my_label}
\end{table}



\subsubsection{Estrutura}
\label{subSub:estrutura}

Realizando um comaprativo entre a elaboração da alu, primeio o modelo mais simples onde as opçoes da ULA são definidas por um simples switch case onde cada saida é separada como no códgio X onde é implemtado uma alu de demostração. Esse modelo de implemtação é visto em diversos trbalhos, como os trabalhos de \citeonline{exem_anbalagan2026rv32i}, \citeonline{exem_hari545543_2026_rv32i}, \citeonline{exem_avelar2026granderisco5} e
\citeonline{exem_calcada2026rvx}.

Para realizar o teste foi elaborado o modulo  superior simples com registradores na entrada e saída para ser posivel utilizar a feramtena de analise de tempo.


Seguindo o mesmo padrão foi realizado um teste elaborando a ALU inpirado na implementação realizada no livro \citeonline{vahid2008sistemas}, nesse modelo é precisso dividir a ALU em dois blocos, um só responsavel pela soma, e outro chamado de extensão, querealiza qualqiuer manipulação nessecesaria. Portanto foi produzido os tres modulos, codigo x contendo sómente a intencia dos dois blocos, o código x contendo a descrição da extensão onde é realizado as modificaçoes, e por fim o codigo x soemnte realizando a soma dos valores e dando o resultado.





Alem dos dois modelos foi explorado o usso de Intenligencia Artificial generativa para tentar obter um resultado com maior velociade, par isso foi solictado as IA Gemini, ChatGpt e Claude para gerarem uma ALU similar as anterios porem do melhor meio para reduzir o tempo



A Tabela \ref{tab:alu_comparacao} foi elaborado com basse nos resultados a analise de implemtação do vivado sendo analisado o Tempo do caminho critino em ns, o consumo total e a quantidade de LUT ocupada. O consumo de energia de energia dos 5 metodos ficaram com valores praticametne iguais, de forma simila rocore nas LUT só tendo duas execoes que abaramutilizando muito mais espaço que os demais, agora o caminho critico sim ha umadifereça bem significativa  einteresante, onde o memtodo mais simples acabou tendo um tempo muito maior queo o nescessario, ja o memtodo do livro ja alcançou um resultado satisfatório, porem ainda é posivel conseguir um pouco mais de velocidade utilizando o modelo gerado pelo gemini, onde conseguiu o menor dos tempo.



\begin{table}[!htb]
  \centering
  \caption{Comparação de Métodos de Implementação de ALU em SystemVerilog no Vivado}
  \label{tab:alu_comparacao}
  \begin{tabular}{l c c c  }
    \toprule
    \textbf{Método / Origem} & \textbf{Caminho Crítico} & \textbf{Potência Total}  & \textbf{LUTs} \\
                             & (ns)                    & (W)                      &               \\
    \midrule
    Método do Livro          & 9,551                   & 0,125                   & 375            \\
    Claude                   & 9,192                   & 0,11                    & 403             \\
    Método Simples           & 7,760                   & 0,12                    & 423            \\
    GPT                      & 7,602                   & 0,111                   & 435            \\
    Gemini                   & 6,839                  & 0,111                   & 357            \\
    \bottomrule
  \end{tabular}
\end{table}


\subsubsection{Operações e I/O}
\label{subSub:operacoes}

Fazendo uma analise das intrunçoes foi determinado as operaçoes que a ULA será capaz de realizar, essas opçoes foram definidasd pensando em melhorar a velocidade das operaçoes, contemplando todas as operaçoes nscessarias confome a Tabela \ref{tab:alu_rv32i}, as operaçoes \textit{ECALL, EBREAK} e \textit{FENCE} não utilizan a ALU toda sua manipulação é relizado diretamente pela unidade de controle. 

Duas intruçoes que foram adicoandas nessa imp,emtação para reduzir ocamnho critico foram a PAS\_B evitando assi m a nescessidade de um mux para realziar a operação LUI sem qualquer interferencia do valor que tiver em A assi removendo a nescessidade de um mux apra garir valor. Outra intrução incomun é a de igualdade para assim conseguir reduzir o caminho critico evitando que todas as saidas da ALU tivesse que garantir esta zerada para sair o sinal de 1 na flag de zero, co mesta operaçã oela faz  ocomaprativo direto e colcoa no baramenteo de saida o resultado, onde o primeiro bit menso significado contem o resultado da operação, esse memso bit é o direciaonado para flag zero, apesar que ao fazer isso em outras intruçoes o valro dessa flag pode estar errado, porem como só para duas intruçoes esse dinal vai importar habilitando fica a cargo da unidade de contrle utilizar ou não. 


\begin{table}[!htb]
\centering
\caption{Mapeamento entre as Operações da ALU e as Instruções do RV32I}
\label{tab:alu_rv32i}
\begin{tabular}{lllp{5cm}}
\toprule
\textbf{Descrição / Função} & \textbf{Operação da ALU} & \textbf{Código Binário} & \textbf{Instruções RISC-V} \\
\midrule
Soma & \texttt{ADD} & \texttt{0000} & 
    \begin{tabular}[t]{@{}l@{}}\texttt{ADD}, \texttt{ADDI}, \texttt{LW}, \texttt{SW}, \texttt{LB}, \texttt{LH}, \\ \texttt{LBU}, \texttt{LHU}, \texttt{SB}, \texttt{SH}, \texttt{AUIPC}\end{tabular} \\
\addlinespace
Subtração & \texttt{SUB} & \texttt{0001} & 
    \texttt{SUB}\\
\addlinespace
E Lógico & \texttt{AND} & \texttt{0010} & 
    \texttt{AND}, \texttt{ANDI} \\
\addlinespace
OU Lógico & \texttt{OR} & \texttt{0011} & 
    \texttt{OR}, \texttt{ORI} \\
\addlinespace
OU Exclusivo Lógico & \texttt{XOR} & \texttt{0100} & 
    \texttt{XOR}, \texttt{XORI} \\
\addlinespace
Deslocamento Esquerdo Lógico & \texttt{SLL} & \texttt{0101} & 
    \texttt{SLL}, \texttt{SLLI} \\
\addlinespace
Deslocamento Direito Lógico & \texttt{SRL} & \texttt{0110} & 
    \texttt{SRL}, \texttt{SRLI} \\
\addlinespace
Deslocamento Direito Aritmético & \texttt{SRA} & \texttt{0111} & 
    \texttt{SRA}, \texttt{SRAI} \\
\addlinespace
Comparação Menor Que (Com Sinal) & \texttt{SLT} & \texttt{1000} & 
    \texttt{SLT}, \texttt{SLTI}, \texttt{BLT}, \texttt{BGE} \\
\addlinespace
Comparação Menor Que (Sem Sinal) & \texttt{SLTU} & \texttt{1001} & 
    \texttt{SLTU}, \texttt{SLTIU}, \texttt{BLTU}, \texttt{BGEU} \\
 \addlinespace
Comparação de igualdade&\texttt{CMP}& \texttt{1011}&\texttt{BEQ}, \texttt{BNE}\\
\addlinespace
Passagem Direta do Operando B & \texttt{PASS\_B} & \texttt{1010} & 
    \texttt{LUI},  \texttt{JAL}, \texttt{JALR} \\

\bottomrule

\end{tabular}
\end{table}


Para este projeto as entradas composta por dois baramenteos de dados  A e B com 32 bits e um baramento de seleção com quatros bits. A saida será composta por um barramento de 32 bits com os dados, e um bit indicando se o resultado da ALU é zerro, se dá a nescessidade desse bit zerro para as intruçoes de desvio \textit{BEQ} e \textit{BNE} onde é precisso realizar uma comparação para ver se os valores são iguais, para isso realiza uma subtração e verifica se o valor é zero, onde só será zero se os dosi valores de entrada forem ezatamente igual, até mesmo o sinal. 

\subsubsection{Validação ALU}
\label{subSub:tb_alu}

para validar a alu finalizada e garatir sua integridade é preciso pasala por um testbench, nesse testbench autmatico é ingetado valores na ula conehcido e comparado com esperados, alem de gerar mil valores aleatório para que memos com varios continua valido os resultados, para isso utilizace uma função dentro do ambiente de simulaçao que tambme recebe as entrada s epoduz uma saida. Ao final o sistema monta u mrelatório infomando se algu mvalor não corepondeu coretamente ao esperado.

inserido onze teste de forma manual  e mil automaticos, todos pasaram com sucesso, valdiando assim o funcionamento da alu.



\subsection{BUSCA}
\label{sub:busca}

Uma intrução de 32 bits é dividida em 4 bytes, onde cada uma fica seguido após a outra como uma pilha de intrução, o é um simples registrador de 32 bits que armazena o endereço que aintrução atual esta dentro da memoria de intrução.

image

o modulo de busca de intrução é relativamente simplples, compsoto por um registrador de progrma (PC, do ingles \textit{Progrman count}), esse registrador guarda o valor da intrução que será 

O tamanho da memoria de intruçoes pode chegar até (valor maximo), porem para esta aplicação não é nescessario tão grande portanto será utilizado apenas 1024 endereços.

para o modulo de busca ficou com cinco entradas e tres saidas. Uma entrada para   o \textit{clock} (clk), uma para reiniciar o modulo (rest_i), uma para permitir a escrita no registrador pc (pq_write_i), essas entradas não serão utilzadas normalmente, são apenas para iniciação e controle manual, um baramento com novo endereço do desvio (new_pc_i) e um bit de controle para realizar o desvio para o endereço do new_pc_i (pc_branch_i), essas etradas é para realizar todas as intruçoes de desvio coretamente.

As saidas serão composta pela intrção que foi lida da memoria (instr_o), o atual valor do pc, que será utilizado para o calcular do endereço de desvio (pc_o) e o valor de pc mais quatro, que representa  oproximo endereço casso não acontença nenhum desvio, esse valor será armazenados nas intruçoes JAL e JAR

O modulo possui funcionanmenteo simples, a cada mudança de clock ele atualiza o valor do contandor do pc, ou para o valor pc + 4 ou para o endereço do desvio. O valor do pc é diividido em tres caminhos, para ler o valro da INTRUÇÃO, soamr 4 ou para saida.

PAra realizar a valição do blocoo foi elaborado um testbanch, onde faço ele ter um comportamento conehcido, que posso prever o resultado e inserir o valor esperado, assim o modulo ja compara o valor do resultado com o valor informado como esperado, como o modulo atualiza seu valor na borda de subida do clock, todas as atualizaçoes no ambiente de simulação devem ocrer no meio, para que assim todos os sinasi que precisarem consigam se propragar coretamente, para isso pode ser utilizado no teste a borda de descida do clock para atualizar os valores e verifica se a alteração que ocoreu na operaçao anterior está coreta.

Todos os teste pasaram perfeitamente nessa estapa.


\subsection{DECODIFICAÇÂO}
\label{sub:decodificar}


\subsection{Execução}
\label{sub:executar}

A unidade de Execução é compsota por duas partes, a ALU e bloco de desvio.

Para comteplar todas as intruções do conjuunto rv32i é preciso que as entradas da ALU alterem dependnendo da intrução, na entrada A para a maioria das intruçoes entra o valor do registrador um, porem para a intrução AUIPC é preciso ter o valor do pc, portanto na entrada A entrar´o pc ou imm, sendo definidio pelo sinal de controle control_alu_a vindo da CU.

A entrda B da ALU, ja pode assumir tres valores, o valrodireto do resgistrador 2 para intruçoes do tipo R e tipo B, para as intruçoes de desvio incondicional será posivel etrarr o valor de pc + 4, para assim registrar ela nos registradores, essa escolha foi tomada com o bjetivco de reduzir o caminho critico evitando um mux na entrada da memoria do registrador, para as demais intruçoes o valor sera o imm deslocado um bit para esquerda, esse deslocamneto é definido pela propria rv322i para garantir o alinhamento da memoria.

a segunda parte do bloco de execução é o gerador de endereço de desvio, esse bolco conta cum um somador proprio, como todas as operaçoes de desvios foram colocadas para esse somador é preciso colocar um mux em uma de sus entradas, de todas as intruçoes de desvio, condicionais ou incondicionais, só uma realiza uma operação diferente de pc + imm, que seria  intrução jalr onde faz rs1 + imm, portanto sóemnte quando for essa intrução um das entradas do somador deve ser trocada para o valor de pc para o registrador 1. 


alem de determinar o endereço para qual sera desviado, esse bloco tambem ´utilizado para determinar se deve ou não realizar o desvio, princiapalmete para os desvios incondicionais, onde esse bloco utiliza o valor da saida zerro da ULA junto com sisnais de controle do CU, para determinar em cada uma das quatro situaçoes, (salto condiconal condição verdaderria, salto coniconal condição falsa, salto incondicioanl relativo ao pc, salto incondicional relativo ao rs1) se deve saltar ou não.



para validar o funcionamento foi realizado um testbench contendo 34 teste, testando todas as operaçoes da ALU, com as diferentes entradas possiveis, e testado os 4 diferentes saltos posiveis, os teste foam ralizados de forma automatica e todos passaram na validação.

\subsection{Memoria de dados}
\label{sub:memoeria_dados}

para a memoria de dados é preciso de seis entadas, uma de clock (clk), uma de controle para habilitar escrita (men_write_i), uma com o tamanho que sera escrito ou lido , 8 bits, 16 bits, ou 32 bits (mem_width_i), uma indicar se o resultado lido deve ser estendido com sinal ou não (mem_signed_i), uam que contennha o endereço de escrita e leitura (addr_i) e uma com os dados que serão gravados (data_i), para saida  s´oé preciso de uma com o resultado.

alem dessas entradas este modulo tambme da suporte para  carega um arqquivo da memoria ja com alguns dados atraves do apremtro MEN_FILE que porpadraõa é data.hex.

O funcioanemtemnteo desse bloco é simpels posuindo um bloco de BRAM da FPGA que fica reponsanvei por realemte armazenar os dados, esse bloco é montado como uma matrix bidireicional, contendo erm cad espaço 8 bits ou um byte, e em cada linha 4 colunas que juntos formam os 32 bits de alinhamento definifido pela RISC-V.

Para relaizar a leitura dos dados é utilziado dos sinais de controle par semper posicionar os bits no primeiros bits e o resto compeltar de acordo com o sian lde extensão. Para escrita é aplicado uam logica simialr onde em cada coluna é alterada individualemte para assim quando for as intruçeos de BYte os dados  ao lado não são alterados. 

Foi elaborado u mtestbench pra esse modul oque vai gravando os dados na memoria e verifica se eles foram gravados coretaemtne, alem de gratir que os as gravaçoe sde Bytes e meia palvras não afetem mais que o nescessario. OS teste foram realizados de forma automatica aplicand as mudanção na descida do clok e verificando se esta tudo certo na subida do clock seguinte, Todos tes passaram.

\subsection{Escrita}
\label{sub:escrita}

O bloco de escita no registradores é o bloco responsavel por escolher qual dado deve ser salvo nos registradores se é um valor vindo da memoria ou um valor vindo do ALU, portato ese modulo só precisa das entrdasm o valor vindo da memoria (mem_data_i), a entrada vindo da saida da ALU (alu_result_i) e uma entrada de controle vindo do bloco de controle (mem_to_reg_i), e para saidas bastauma saida com a escolha do resutlado (wb_data_o).

Quando o valor da entrada mem_to_reg_i estiver 0 a saida sera ovalor do resultado da alu, e quanfo for 1 saira o valro mem_to_reg_i. 

Algumas implemtnaçoe  observadas, como a do \citeonline{exem_hari545543_2026_rv32i}e do \citeonline{exem_calcada2026rvx} utilizan nessa etapa tambem para entrada do valor com contador de programa mais quatro, porem com ofoi adiconado direto na entrada do bloco de execução com o bejtivo de otimizar os seletores não será preciso aqui.

Para validação desse bloco foi elaborado um testebench simple automatizado aplicando mudnaa do valor na borda de decida e ja verificando o resultado. todos teste passaram.

\subsection{Decodificação}
\label{sub:decodificacao}


\subsubsection{Registrador}
\label{subsub:registrador}

Para o modulo de registro ficoiu com uma entrada para \textit{clock} (clk), uma entrada parahabilitar a escrita (written_enable_i), um baramenteo de 32 bits com os dados a serem escritos (written_data_i), um abarementeo com 5 bits por padrão com o endereço de escrita (written_add_i), outros dois baramento de 5 bits co mendereço do registro que ddeve ser lido no rs1 e rs2 (rs1_add_i e rs2_add_i) e apra saida doi baramento de 32 bits com os dados lidos (rs1_o e rs2_o).

Este modulo foi elaborado de form asimpels utilziando a declaração de como array de 31 posiçoes (1:31) ja que o registradr 0 por definição da RV31I deve ser semore 0 sendo imposivel de escrever nele, assim com o bjetivo de simplificar foi gerado apenas 31 valores.

Tanto para escrita quandot leitura foi adicionado um verificador de condição para garantir que quando o endereço seja 0 de 0 na saida e não faça nada na entrada, apesar de adiianr um bloco de logica mais alem da propria seleção antural da estrutura de array após compilação do código o resultado é otimizado e esse bloco desaparece se tornand parte do conjunto compelto de registradores.

Para fim de garatir foi realizado um pequeno testbench automatico que valida se os dados estão sendo salvos coretamente no registrador e as leituras então coretas, todas a lteraçoes de siansi raealizadas na bordsa de descida e todas as verificaçoes na borda de subiida, ja que este é um compoennte sincrino de borda de subida. 

\subsubsection{Gerado de Imaediato}
\label{subsub:imm_gen}


O gerador de imediato é o bloco capaz de pegar um valor que está no meio da intrução e deixalo formato e com tamanho coreto para execução para o resto do circuito. Para esse gerador é preciso somente entrda da intrução completa com 32 bits (instruction_i), e a saida com 32 bits do valor pronto (imm_o).

O formato da decodificaçã ode uma intrução varia confome o tipo e variação de ada uma, e para identificar o tipo basta a nalisar o opconde da intrução (7 priemrios bits) e assim codificar ela para o capçar confome especificado pela isa rv32i.

Foi realizado um testbench automatico para este modulo onde os dados são inseridos na borda de escida e ja verificado a resposta, todos os teste passaram coretamente.

\subsubsection{Control unit}
\label{subsub:control_unit}



A unidade de controle (\textit{control unit}) recebe a instrução completa de 32 bits por meio da entrada \texttt{instruction\_i}. A partir dos campos \textit{opcode}, \textit{func3} e \textit{func7}, o bloco identifica a operação e gera os sinais usados pelas etapas de execução, acesso à memória, escrita no banco de registradores e atualização do contador de programa. A unidade é combinacional e não possui entrada de clock.

O formato geral da instrução e os campos consultados pela unidade são apresentados na Tabela~\ref{tab:control-unit-fields}. A interpretação dos demais campos depende do formato da instrução: em instruções do tipo I, S, B, U e J, parte dos bits que ocupa a posição de \textit{func7} em instruções do tipo R pertence ao imediato.

\begin{table}[!htb]
\centering
\caption{Campos da instrução consultados pela unidade de controle.}
\label{tab:control-unit-fields}
\begin{tabular}{lll}
\toprule
\textbf{Campo} & \textbf{Bits} & \textbf{Uso na decodificação} \\
\midrule
\texttt{opcode} & \texttt{[6:0]} & Identifica a classe da instrução. \\
\texttt{func3} & \texttt{[14:12]} & Distingue operações dentro de uma classe. \\
\texttt{func7} & \texttt{[31:25]} & Distingue ADD/SUB, SRL/SRA e SRLI/SRAI. \\
\bottomrule
\end{tabular}
\end{table}

As portas da unidade e sua relação com os blocos do processador estão descritas na Tabela~\ref{tab:control-unit-ports}. Os prefixos dos nomes das saídas indicam o módulo que recebe cada sinal. Os sinais de largura e extensão de memória usam os tipos definidos em \texttt{mem\_pkg}; os sinais da ALU e dos desvios usam os tipos definidos em \texttt{alu\_pkg} e \texttt{branch\_pkg}.

\begin{table}[!htb]
\centering
\caption{Entradas e saídas da unidade de controle.}
\label{tab:control-unit-ports}
\begin{tabular}{p{4.1cm}p{2.1cm}p{8.0cm}}
\toprule
\textbf{Sinal} & \textbf{Direção / largura} & \textbf{Descrição} \\
\midrule
\texttt{instruction\_i} & Entrada, 32 bits & Instrução completa a ser decodificada. \\
\texttt{execution\_alu\_control\_o} & Saída, \texttt{alu\_op\_t} & Seleciona a operação da ALU. \\
\texttt{execution\_branch\_control\_o} & Saída, \texttt{branch\_op\_t} & Seleciona a condição ou a origem do desvio. \\
\texttt{execution\_branch\_enable\_o} & Saída, 1 bit & Habilita a avaliação de desvio em \texttt{execution}. \\
\texttt{execution\_control\_alu\_a\_o} & Saída, 1 bit & Seleciona rs1 (0) ou PC (1) para a entrada A da ALU. \\
\texttt{execution\_control\_alu\_b\_o} & Saída, 2 bits & Seleciona rs2 (\texttt{00}), imediato (\texttt{01}) ou PC+4 (\texttt{10}) para a entrada B da ALU. \\
\texttt{search\_pc\_write\_o} & Saída, 1 bit & Habilita a atualização do PC; nesta implementação permanece em 1. \\
\texttt{data\_memory\_mem\_write\_o} & Saída, 1 bit & Habilita a escrita na memória de dados. \\
\texttt{data\_memory\_mem\_width\_o} & Saída, \texttt{mem\_width\_t} & Define acesso a byte, meia palavra ou palavra. \\
\texttt{data\_memory\_mem\_signed\_o} & Saída, \texttt{mem\_signed\_t} & Define extensão com ou sem sinal durante loads. \\
\texttt{register\_file\_written\_enable\_o} & Saída, 1 bit & Habilita a escrita do resultado no registrador destino. \\
\texttt{write\_back\_mem\_to\_reg\_o} & Saída, 1 bit & Seleciona dado da memória (1) ou resultado da ALU (0) no write-back. \\
\bottomrule
\end{tabular}
\end{table}

Por padrão, a unidade seleciona \texttt{OP\_ADD}, mantém os desvios desabilitados, não escreve na memória nem no banco de registradores e seleciona o resultado da ALU para write-back. O seletor de largura de memória assume \texttt{MEM\_WIDTH\_WORD}, e o seletor de extensão assume \texttt{MEM\_UNSIGNED}. O sinal \texttt{search\_pc\_write\_o} permanece ativo. Esses valores também são mantidos para opcodes não suportados.

Nas tabelas de decodificação, os campos \textit{opcode}, \textit{func3} e \textit{func7} mostram a codificação da instrução RV32I. Um traço (--) indica que o campo não é aplicável àquele formato ou não participa da seleção descrita. Salvo quando indicado, as saídas não mencionadas em cada classe mantêm os valores padrão apresentados acima.

\paragraph{Operações da ALU.}
A Tabela~\ref{tab:control-unit-alu-decode} resume as instruções aritméticas, lógicas e de deslocamento. Em todas elas, a escrita no registrador destino é habilitada. Para instruções do tipo R, a entrada A recebe rs1 e a entrada B recebe rs2. Para instruções do tipo I, A recebe rs1 e B recebe o imediato. O campo \textit{func7} só diferencia as operações indicadas na tabela; para as demais operações, a implementação não usa esse campo na seleção da ALU.

\begin{table}[!htb]
\centering
\caption{Decodificação das instruções que utilizam a ALU.}
\label{tab:control-unit-alu-decode}
\begin{tabular}{lllll}
\toprule
\textbf{Instrução} & \textbf{opcode} & \textbf{func3} & \textbf{func7} & \textbf{Operação da ALU} \\
\midrule
ADD   & \texttt{0110011} & \texttt{000} & \texttt{0000000} & \texttt{OP\_ADD} \\
SUB   & \texttt{0110011} & \texttt{000} & \texttt{0100000} & \texttt{OP\_SUB} \\
SLL   & \texttt{0110011} & \texttt{001} & \texttt{0000000} & \texttt{OP\_SLL} \\
SLT   & \texttt{0110011} & \texttt{010} & \texttt{0000000} & \texttt{OP\_SLT} \\
SLTU  & \texttt{0110011} & \texttt{011} & \texttt{0000000} & \texttt{OP\_SLTU} \\
XOR   & \texttt{0110011} & \texttt{100} & \texttt{0000000} & \texttt{OP\_XOR} \\
SRL   & \texttt{0110011} & \texttt{101} & \texttt{0000000} & \texttt{OP\_SRL} \\
SRA   & \texttt{0110011} & \texttt{101} & \texttt{0100000} & \texttt{OP\_SRA} \\
OR    & \texttt{0110011} & \texttt{110} & \texttt{0000000} & \texttt{OP\_OR} \\
AND   & \texttt{0110011} & \texttt{111} & \texttt{0000000} & \texttt{OP\_AND} \\
\addlinespace
ADDI  & \texttt{0010011} & \texttt{000} & -- & \texttt{OP\_ADD} \\
SLTI  & \texttt{0010011} & \texttt{010} & -- & \texttt{OP\_SLT} \\
SLTIU & \texttt{0010011} & \texttt{011} & -- & \texttt{OP\_SLTU} \\
XORI  & \texttt{0010011} & \texttt{100} & -- & \texttt{OP\_XOR} \\
ORI   & \texttt{0010011} & \texttt{110} & -- & \texttt{OP\_OR} \\
ANDI  & \texttt{0010011} & \texttt{111} & -- & \texttt{OP\_AND} \\
SLLI  & \texttt{0010011} & \texttt{001} & \texttt{0000000} & \texttt{OP\_SLL} \\
SRLI  & \texttt{0010011} & \texttt{101} & \texttt{0000000} & \texttt{OP\_SRL} \\
SRAI  & \texttt{0010011} & \texttt{101} & \texttt{0100000} & \texttt{OP\_SRA} \\
\bottomrule
\end{tabular}
\end{table}

Nas linhas em que \textit{func7} aparece como --, esse campo não é verificado pela unidade para selecionar a operação. Em SLLI e nos deslocamentos à direita imediatos, os bits superiores do imediato ocupam a posição correspondente a \textit{func7}; a unidade compara esses bits com os padrões de deslocamento definidos em \texttt{func7\_pkg}.

\paragraph{Acessos à memória.}
Para loads e stores, a ALU soma rs1 ao imediato, com \texttt{execution\_control\_alu\_b\_o = 01}. A Tabela~\ref{tab:control-unit-memory-decode} apresenta os campos e sinais específicos de cada operação.

% Tabela de acessos à memória: trecho não incluído no texto recebido.

